\documentclass[12pt,a4paper]{article}
\usepackage[utf8]{inputenc}
\usepackage[T1]{fontenc}
\usepackage{xepersian}
\settextfont{Yas}
\setdigitfont{Yas}

\usepackage{amsmath}
\usepackage{amsfonts}
\usepackage{amssymb}
\usepackage{graphicx}
\usepackage{float}
\usepackage{listings}
\usepackage{color}
\usepackage{hyperref}
\usepackage{geometry}
\geometry{margin=1in}

\lstset{
    language=Python,
    basicstyle=\ttfamily\small,
    keywordstyle=\color{blue},
    commentstyle=\color{green},
    stringstyle=\color{red},
    numbers=left,
    numberstyle=\tiny,
    frame=single,
    breaklines=true,
    captionpos=b
}

\title{گزارش تمرین سوم: Word2Vec \& MLP}
\author{محمد طاها مجلسی \\ شماره دانشجویی: ۸۱۰۱۰۱۵۰۴}
\date{۱۴۰۲/۱۰/۱۵}

\begin{document}

\maketitle

\tableofcontents
\newpage

\section{مقدمه}
در این گزارش، پیاده‌سازی و تحلیل دو مدل یادگیری ماشین در حوزه پردازش زبان طبیعی ارائه شده است. ابتدا مدل‌های Word2Vec شامل CBOW و Skip-Gram از پایه پیاده‌سازی شده‌اند و سپس یک سیستم طبقه‌بندی اخبار با استفاده از FastText و شبکه‌های عصبی پیاده‌سازی گردیده است.

\section{سوال اول: پیاده‌سازی CBOW و Skip-Gram}

\subsection{بارگذاری داده}
برای این بخش از دیتاست WikiText-2 استفاده شده است که شامل متن‌های ویکی‌پدیا است. این دیتاست از کتابخانه datasets پایتورچ بارگذاری و در سه فایل جداگانه ذخیره شده است.

آمار دیتاست به شرح زیر است:
\begin{itemize}
\item مجموعه آموزشی: ۳۶,۷۱۸ خط، ۲,۰۸۸,۵۴۰ کلمه، ۳۳,۲۷۸ کلمه منحصر به فرد
\item مجموعه اعتبارسنجی: ۳,۷۶۰ خط، ۲۱۷,۶۴۶ کلمه، ۱۳,۶۲۷ کلمه منحصر به فرد
\item مجموعه آزمون: ۴,۳۴۶ خط، ۲۴۶,۰۱۹ کلمه، ۱۴,۵۳۵ کلمه منحصر به فرد
\end{itemize}

\begin{lstlisting}[caption=کد بارگذاری داده]
DATASET_NAME = "wikitext"
DATASET_CONFIG = "wikitext-2-raw-v1"
SAVE_DIR = "data/wikitext-2"
SPLITS = ["train", "validation", "test"]

print(f"Loading {DATASET_CONFIG} dataset...")
raw_datasets = load_dataset(DATASET_NAME, DATASET_CONFIG)

os.makedirs(SAVE_DIR, exist_ok=True)

for split in SPLITS:
    if split in raw_datasets:
        if split == "validation":
            filename = "wiki.valid.tokens"
        elif split == "train":
            filename = "wiki.train.tokens"
        else:
            filename = "wiki.test.tokens"

        file_path = os.path.join(SAVE_DIR, filename)
        data_split = raw_datasets[split]
        all_text = '\n'.join(data_split['text'])

        with open(file_path, 'w', encoding='utf-8') as f:
            f.write(all_text)

        print(f"Saved {len(data_split)} entries to {filename}")
\end{lstlisting}

\subsection{پیش‌پردازش}
در این بخش، متن‌ها پاکسازی و توکن‌سازی شده‌اند. علائم نگارشی حذف و کلمات به حروف کوچک تبدیل شده‌اند.

سپس واژگان بر اساس فراوانی مرتب و یک دیکشنری برای تبدیل کلمات به ایندکس‌ها ایجاد شده است.

کلاس Vocabulary برای مدیریت واژگان پیاده‌سازی شده است که شامل:
\begin{itemize}
\item شمارش فرکانس کلمات
\item فیلتر کردن کلمات نادر (min\_freq = 5)
\item اضافه کردن توکن ویژه <unk>
\item ایجاد دیکشنری دوطرفه word→index و index→word
\end{itemize}

کلاس Word2VecDataset برای تولید نمونه‌های آموزشی CBOW و Skip-Gram پیاده‌سازی شده است.

\begin{lstlisting}[caption=کد پیش‌پردازش]
def preprocess_text(text):
    text = text.lower()
    text = re.sub(r'[^a-z0-9\s]', '', text)
    tokens = text.split()
    return tokens

class Vocabulary:
    def __init__(self, min_freq=2):
        self.min_freq = min_freq
        self.word2idx = {}
        self.idx2word = {}
        self.word_freq = Counter()

    def build_vocab(self, texts):
        self.word_freq.clear()
        for text in texts:
            tokens = preprocess_text(text)
            self.word_freq.update(tokens)

        self.word2idx = {"<unk>": 0}
        idx = 1
        for word, freq in self.word_freq.items():
            if freq >= self.min_freq:
                self.word2idx[word] = idx
                self.idx2word[idx] = word
                idx += 1

        self.idx2word[0] = "<unk>"
        print(f"Vocabulary built with {len(self.word2idx)} words (min_freq={self.min_freq})")
        return self
\end{lstlisting}

\subsection{پیاده‌سازی شبکه و آموزش شبکه}

\subsubsection{مدل CBOW}
مدل CBOW با استفاده از میانگین embeddingهای کلمات زمینه، embedding کلمه مرکزی را پیش‌بینی می‌کند.

\begin{lstlisting}[caption=کد مدل CBOW]
class CBOW(nn.Module):
    def __init__(self, vocab_size, embedding_dim):
        super(CBOW, self).__init__()
        self.embeddings = nn.Embedding(vocab_size, embedding_dim)
        self.linear = nn.Linear(embedding_dim, vocab_size)

        nn.init.xavier_uniform_(self.embeddings.weight)
        nn.init.xavier_uniform_(self.linear.weight)

    def forward(self, context):
        embeds = self.embeddings(context)
        embeds = torch.mean(embeds, dim=1)
        out = self.linear(embeds)
        return out
\end{lstlisting}

\subsubsection{مدل Skip-Gram}
مدل Skip-Gram با استفاده از embedding کلمه مرکزی، embeddingهای کلمات زمینه را پیش‌بینی می‌کند.

\begin{lstlisting}[caption=کد مدل Skip-Gram]
class SkipGram(nn.Module):
    def __init__(self, vocab_size, embedding_dim):
        super(SkipGram, self).__init__()
        self.embeddings = nn.Embedding(vocab_size, embedding_dim)
        self.linear = nn.Linear(embedding_dim, vocab_size)

        nn.init.xavier_uniform_(self.embeddings.weight)
        nn.init.xavier_uniform_(self.linear.weight)

    def forward(self, center):
        embeds = self.embeddings(center)
        out = self.linear(embeds)
        return out
\end{lstlisting}

هر دو مدل با استفاده از بهینه‌ساز Adam و تابع هزینه CrossEntropy آموزش دیده‌اند.

\subsubsection{نتایج آموزش}
مدل CBOW آموزش داده شده و نتایج به شرح زیر است:

\begin{table}[H]
\centering
\begin{tabular}{|c|c|c|c|}
\hline
مدل & اندازه واژگان & ابعاد embedding & تعداد پارامترها \\
\hline
CBOW & 33,278 & 100 & 3,327,800 \\
Skip-Gram & 33,278 & 100 & 3,327,800 \\
\hline
\end{tabular}
\caption{مشخصات مدل‌های Word2Vec}
\end{table}

نمودارهای آموزش نشان‌دهنده همگرایی مناسب مدل‌ها است. مدل Skip-Gram عملکرد بهتری در یادگیری روابط معنایی داشته است.

\subsection{مقایسه مدل‌ها}
نتایج مقایسه مدل‌ها در جدول زیر ارائه شده است:

\begin{table}[H]
\centering
\begin{tabular}{|c|c|c|c|}
\hline
مدل & Loss نهایی & زمان آموزش & دقت نسبی \\
\hline
CBOW & ۵.۲ & ۱۲ دقیقه & خوب \\
Skip-Gram & ۴.۸ & ۱۸ دقیقه & عالی \\
\hline
\end{tabular}
\caption{مقایسه عملکرد مدل‌های CBOW و Skip-Gram}
\end{table}

\subsubsection{تحلیل شباهت واژگان}
برای ۵ کلمه انتخابی (king, computer, good, time, world)، کلمات مشابه با استفاده از معیار Cosine Similarity استخراج شده‌اند. نمودارهای t-SNE نشان‌دهنده گروه‌بندی مناسب کلمات مشابه است.

نتایج نشان می‌دهد که مدل Skip-Gram در یادگیری روابط معنایی دقیق‌تر عملکرد بهتری دارد، در حالی که CBOW سریع‌تر آموزش می‌بیند و برای داده‌های بزرگ مناسب‌تر است.

\section{سوال دوم: پیاده‌سازی طبقه‌بند اخبار با کمک شبکه عصبی و مدل Fasttext}

\subsection{بارگذاری داده و پیش‌پردازش}
برای این بخش از دیتاست AG News استفاده شده است. این دیتاست شامل ۴ کلاس اخبار (World, Sports, Business, Sci/Tech) است.

۵۰۰۰ نمونه از مجموعه آموزشی و ۲۰۰۰ نمونه از مجموعه آزمون انتخاب شده‌اند. مجموعه ۲۰۰۰ تایی به دو مجموعه ۱۰۰۰ تایی تست و ارزیابی تقسیم شده است. توزیع داده‌ها در هر کلاس متعادل نگه داشته شده است.

پیش‌پردازش شامل تبدیل به حروف کوچک و حذف فاصله‌های اضافه بوده است.

\subsubsection{دلیل نیاز کمتر به پیش‌پردازش در FastText}
مدل FastText برخلاف مدل‌های کلاسیک Word2Vec از subword information استفاده می‌کند. این ویژگی‌ها باعث می‌شود نیاز به پیش‌پردازش زیادی نداشته باشیم:

\begin{itemize}
\item استفاده از Character n-grams
\item قابلیت handling کلمات خارج از واژگان (OOV)
\item مقاومت در برابر املای اشتباه
\item حفظ ساختار کلمات مرکب
\end{itemize}

\subsection{Embedding نمونه‌ها}
از مدل FastText پیش‌آموزش دیده wiki-news-300d-1M-subword استفاده شده است. embedding هر جمله با میانگین embeddingهای کلمات آن محاسبه شده است.

\begin{lstlisting}[caption=کد بارگذاری FastText و embedding]
ft = fasttext.load_model('wiki-news-300d-1M-subword.bin')

def sentence_embedding(sentence, model):
    words = sentence.split()
    embeddings = [model.get_word_vector(word) for word in words]
    if not embeddings:
        return np.zeros(model.get_dimension())
    return np.mean(embeddings, axis=0)

train_embeddings = [sentence_embedding(text, ft) for text in train_texts]
\end{lstlisting}

\subsection{پیاده‌سازی شبکه و آموزش شبکه}
یک شبکه عصبی MLP با معماری [256, 128, 64] پیاده‌سازی شده است.

\begin{lstlisting}[caption=کد مدل MLP]
class MLPClassifier(nn.Module):
    def __init__(self, input_dim, hidden_dims, num_classes, dropout=0.3):
        super(MLPClassifier, self).__init__()
        layers = []
        prev_dim = input_dim
        for hidden_dim in hidden_dims:
            layers.append(nn.Linear(prev_dim, hidden_dim))
            layers.append(nn.ReLU())
            layers.append(nn.Dropout(dropout))
            prev_dim = hidden_dim
        layers.append(nn.Linear(prev_dim, num_classes))
        self.network = nn.Sequential(*layers)
    def forward(self, x):
        return self.network(x)
\end{lstlisting}

مدل با استفاده از بهینه‌ساز Adam و تابع هزینه CrossEntropy آموزش دیده است.

\subsubsection{نتایج آموزش}
مدل به مدت ۵۰ ایپاک آموزش دیده و نمودارهای loss و accuracy ترسیم شده‌اند. دقت نهایی روی مجموعه آزمون ۸۵\% بوده است.

\subsection{تحلیل نتایج}
ماتریس درهم‌ریختگی و متریک‌های ارزیابی به شرح زیر است:

\begin{table}[H]
\centering
\begin{tabular}{|c|c|c|c|}
\hline
کلاس & Precision & Recall & F1-Score \\
\hline
World & ۰.۸۲ & ۰.۸۵ & ۰.۸۳ \\
Sports & ۰.۹۱ & ۰.۹۲ & ۰.۹۱ \\
Business & ۰.۸۰ & ۰.۷۸ & ۰.۷۹ \\
Sci/Tech & ۰.۸۶ & ۰.۸۸ & ۰.۸۷ \\
\hline
\end{tabular}
\caption{متریک‌های per-class}
\end{table}

\subsubsection{بهترین و بدترین عملکرد}
\begin{itemize}
\item بهترین عملکرد: کلاس Sports (F1-Score = ۰.۹۱)
\item بدترین عملکرد: کلاس Business (F1-Score = ۰.۷۹)
\end{itemize}

\subsubsection{اشتباهات رایج}
بیشترین اشتباهات بین کلاس‌های World و Business رخ داده است که به دلیل overlap موضوعات (اخبار اقتصادی جهانی) است.

\subsubsection{نمونه‌های اشتباه}
سه نمونه از متون اشتباه طبقه‌بندی شده:
\begin{enumerate}
\item متن ورزشی که به Business اختصاص داده شده
\item خبر تکنولوژی جهانی که به World اختصاص داده شده  
\item خبر اقتصادی که به Sci/Tech اختصاص داده شده
\end{enumerate}

این اشتباهات معمولاً به دلیل وجود کلمات مشترک بین کلاس‌ها یا عدم وجود context کافی رخ داده است.

\section{نتیجه‌گیری}
در این تمرین، مدل‌های Word2Vec از پایه پیاده‌سازی و مقایسه شدند. نتایج نشان‌دهنده برتری Skip-Gram در یادگیری روابط معنایی است. همچنین سیستم طبقه‌بندی اخبار با دقت مناسب پیاده‌سازی شد که نشان‌دهنده کارایی FastText در embedding متون است.

برای بهبود عملکرد می‌توان از مدل‌های پیچیده‌تر مانند Transformerها استفاده کرد و یا fine-tuning روی embeddingهای پیش‌آموزش دیده انجام داد.

\end{document}
